{\actuality} Одним из важных компонентов динамического тренажера автотранспорта и спецтехники (АТ и СТ) является система подвижности (СП), предназначенная для выработки у обучаемого необходимых моторных навыков в процессе управления техникой. Применяемые в большинстве тренажеров системы подвижности представляют собой манипулятор с параллельной кинематикой, состоящий из подвижной платформы, связанной с неподвижным основанием одним или несколькими поступательными звеньями. В зависимости от требований к имитируемым эффектам, используются платформы с числом степеней свободы от 1 до 6. В подавляющем большинстве тренажеров АТ и СТ (порядка 70-80~\% продукции ведущих производителей тренажерных систем) применяются системы подвижности 2DOF и 3DOF с кривошипно-шатунным механизмом (КШМ) привода звеньев, обеспечивающие отработку необходимых при имитации данной техники координат подвижной платформы (крен, тангаж, высота). При разработке ПО управления универсальной системой подвижности (УСП) необходимо учитывать такие факторы, как геометрические параметры платформы, рабочую зону манипулятора, массогабаритные характеристики платформы в различных состояниях, алгоритмы управления положением привода, протоколы обмена с датчиками и исполнительными устройствами. При наличии в номенклатуре производимых тренажеров нескольких систем подвижности с различными кинематическими схемами и типами приводов достаточно остро встает вопрос о разработке ПО, позволяющего осуществлять управление несколькими СП по стандартизованному протоколу обмена и производить быструю настройку под конкретную СП с возможностью диагностики состояния оборудования.	
 
 \textbf{Актуальность} рассматриваемой диссертации определяется потребностями:
\begin{itemize}
	\item создания систем подвижности, обеспечивающих наиболее реалистичные эффекты движения имитируемой техники;
	\item обеспечения безопасной эксплуатации АРМ обучаемого как подвижного объекта с людьми, что требует применения специальных мер (дополнительные датчики, дублирование систем);
	\item унификации модулей тренажерного комплекса, в частности, системы подвижности, что позволяет существенно упростить процесс развертывания тренажера в ходе пуско-наладочных работ.
В работе рассматриваются вопросы реализации системы управления несколькими типами СП, построение архитектуры приложения, отвечающей заданным требованиям универсализации, а также применяемые при этом инженерные решения в области аппаратной части платформы.
\end{itemize}
% {\progress} 
% Этот раздел должен быть отдельным структурным элементом по
% ГОСТ, но он, как правило, включается в описание актуальности
% темы. Нужен он отдельным структурынм элемементом или нет ---
% смотрите другие диссертации вашего совета, скорее всего не нужен.

{\aim} данной работы является разработка унифицированной системы управления динамической платформой подвижности, позволяющей абстрагировать реализацию алгоритмов управления и ограничений, накладываемых рабочей зоной платформы, от клиентской стороны -- модуля тренажерного комплекса, отвечающего за преобразование и передачу данных телеметрии от математической модели техники к системе подвижности

Для~достижения поставленной цели необходимо было решить следующие {\tasks}:
\begin{enumerate}
  \item разработать алгоритмы управления СП, учитывающие динамические характеристики подвижной платформы и закрепленного на ней объекта (соответствует п. 5 паспорта научной специальности);
  \item разработать методику настройки регуляторов СП, обеспечивающие требуемое качество процессов в системе управления (соответствует п. 4 паспорта научной специальности);
  \item разработать алгоритмы формирования траектории в соответствии с рассчитанной рабочей зоной платформы (соответствует п. 6 паспорта научной специальности);
  \item произвести параметрическую идентификацию модели СП по данным от датчиков (соответствует п. 8 паспорта научной специальности);
  \item разработать процедуры автоматической диагностики состояния оборудования СП (соответствует п. 15 паспорта научной специальности);
  \item разработать алгоритмы взаимодействия клиенткой стороны и контроллера движения СП (соответствует п. 17 паспорта научной специальности).
\end{enumerate}

{\novelty}
\begin{enumerate}
  \item получено аналитическое решение уравнений динамики трехстепенной электромеханической платформы с кривошипно-шатунных звеньями и пружинной опорой, пригодное для расчета в режиме реального времени, что позволяет осуществить коррекцию траектории. Это значительно улучшает отклик системы подвижности и уменьшает ошибку регулирования;
  \item разработан алгоритм коррекции траектории движения УСП, учитывающий кинематику и динамику механизма в совокупности с ограничениями приводных систем. Алгоритм предполагает "мягкую" срезку заданных перемещений, скоростей и ускорений, что обеспечивает наилучшее приближение реальной траектории платформы к заданной;
  \item получено аналитическое решение прямой задачи кинематики, позволяющее с достаточной точностью вычислять в режиме реального времени текущие угловые и поступательные координаты СП без применения итерационных алгоритмов, что позволяет не только визуализировать текущее положение платформы без использования дополнительных датчиков, но и применять продвинутые методики управления с обратной связью по координатам платформы;
  \item проведено исследование точности позиционирования СП при помощи прецизионного промышленного датчика углов наклона, экспериментально подтвердившее обоснованность использования описанного выше решения прямой задачи кинематики в задачах управления платформой.
\end{enumerate}

{\influence} определяется внедрением ПАК управления системой подвижности в тренажерных комплексах военной, гражданской техники и речных судов производства АО «Кронштадт Технологии». Результаты работы могут быть использованы при разработке систем управления подвижными платформами тренажеров автотранспорта и спецтехники.

{\methods} Для решения задач, рассматриваемых в работе использовались методы теории систем автоматического управления, линейной алгебры, вычислительной математики. Для создания программно-аппаратного комплекса применены методы технологии программирования. В ходе верификации результатов моделирования использовались методы планирования вычислительных экспериментов и методы обработки данных натурных испытаний.

{\defpositions}
\begin{itemize}
  \item упрощенная модель динамики трехстепенной динамической платформы на кривошипно-шатунных опорах, позволяющая с приемлемой точностью рассчитывать вращающие моменты двигателей по заданным перемещениям, скоростям и ускорениям СП в режиме реального времени; 
  \item алгоритм управления координатами УСП, учитывающий динамику механизма и приводных систем;
  \item алгоритм ограничения координат и безударного выхода на границу рабочей зоны УСП;
  \item методика настройки параметров регулятора с целью достижения заданных требований к УСП;
\end{itemize}

{\reliability} Достоверность исследований, проведенных в диссертации, подтверждается:
\begin{itemize}
	\item проверкой теоретических выкладок и уравнений с помощью наборов входных тестовых данных;
	\item проведенными экспериментальными исследованиями с паспортизацией предельных достижимых характеристик УСП;
	\item тестированием работы УСП специалистами Компании АО «Кронштадт Технологии».
\end{itemize}

{\probation}

Основные результаты диссертационной работы докладывались и обсуждались на конференциях:
\begin{itemize}
	\item на конференции VII Поляховские Чтения – г. Санкт-Петербург, февраль 2015 г.
	\item на конференции Сетевое партнерство в науке, промышленности и образовании – г. Санкт-Петербург, июль 2016 г.,
\end{itemize}

выставках:
\begin{itemize}
	\item Строительная Техника и Технологии (СТТ), 2015, г. Москва,
	\item Международный военно-технический форум «АРМИЯ-2015», 2015 г., г. Кубинка, Московская область,
	\item Ежегодное совещание командования инженерных войск по округам, 2015 г., г. Кстово, Нижегородская область ;
\end{itemize}

\ifnumequal{\value{bibliosel}}{0}
{%%% Встроенная реализация с загрузкой файла через движок bibtex8. (При желании, внутри можно использовать обычные ссылки, наподобие `\cite{vakbib1,vakbib2}`).
    {\publications} Основные результаты по теме диссертации изложены
    в~XX~печатных изданиях,
    X из которых изданы в журналах, рекомендованных ВАК,
    X "--- в тезисах докладов.
}%
{%%% Реализация пакетом biblatex через движок biber
    \begin{refsection}[bl-author, bl-registered]
        % Это refsection=1.
        % Процитированные здесь работы:
        %  * подсчитываются, для автоматического составления фразы "Основные результаты ..."
        %  * попадают в авторскую библиографию, при usefootcite==0 и стиле `\insertbiblioauthor` или `\insertbiblioauthorgrouped`
        %  * нумеруются там в зависимости от порядка команд `\printbibliography` в этом разделе.
        %  * при использовании `\insertbiblioauthorgrouped`, порядок команд `\printbibliography` в нём должен быть тем же (см. biblio/biblatex.tex)
        %
        % Невидимый библиографический список для подсчёта количества публикаций:
        \phantom{\printbibliography[heading=nobibheading, section=1, env=countauthorvak,          keyword=biblioauthorvak]%
        \printbibliography[heading=nobibheading, section=1, env=countauthorwos,          keyword=biblioauthorwos]%
        \printbibliography[heading=nobibheading, section=1, env=countauthorscopus,       keyword=biblioauthorscopus]%
        \printbibliography[heading=nobibheading, section=1, env=countauthorconf,         keyword=biblioauthorconf]%
        \printbibliography[heading=nobibheading, section=1, env=countauthorother,        keyword=biblioauthorother]%
        \printbibliography[heading=nobibheading, section=1, env=countregistered,         keyword=biblioregistered]%
        \printbibliography[heading=nobibheading, section=1, env=countauthorpatent,       keyword=biblioauthorpatent]%
        \printbibliography[heading=nobibheading, section=1, env=countauthorprogram,      keyword=biblioauthorprogram]%
        \printbibliography[heading=nobibheading, section=1, env=countauthor,             keyword=biblioauthor]%
        \printbibliography[heading=nobibheading, section=1, env=countauthorvakscopuswos, filter=vakscopuswos]%
        \printbibliography[heading=nobibheading, section=1, env=countauthorscopuswos,    filter=scopuswos]}%
        %
        \nocite{*}%
        %
        {\publications} Основные результаты по теме диссертации изложены в~\arabic{citeauthor}~печатных изданиях,
        \arabic{citeauthorvak} из которых изданы в журналах, рекомендованных ВАК%
        \ifnum \value{citeauthorscopuswos}>0%
            , \arabic{citeauthorscopuswos} "--- в~периодических научных журналах, индексируемых Web of~Science и Scopus%
        \fi%
        \ifnum \value{citeauthorconf}>0%
            , \arabic{citeauthorconf} "--- в~тезисах докладов.
        \else%
            .
        \fi%
        \ifnum \value{citeregistered}=1%
            \ifnum \value{citeauthorpatent}=1%
                Зарегистрирован \arabic{citeauthorpatent} патент.
            \fi%
            \ifnum \value{citeauthorprogram}=1%
                Зарегистрирована \arabic{citeauthorprogram} программа для ЭВМ.
            \fi%
        \fi%
        \ifnum \value{citeregistered}>1%
            Зарегистрированы\ %
            \ifnum \value{citeauthorpatent}>0%
            \formbytotal{citeauthorpatent}{патент}{}{а}{}%
            \ifnum \value{citeauthorprogram}=0 . \else \ и~\fi%
            \fi%
            \ifnum \value{citeauthorprogram}>0%
            \formbytotal{citeauthorprogram}{программ}{а}{ы}{} для ЭВМ.
            \fi%
        \fi%
        % К публикациям, в которых излагаются основные научные результаты диссертации на соискание учёной
        % степени, в рецензируемых изданиях приравниваются патенты на изобретения, патенты (свидетельства) на
        % полезную модель, патенты на промышленный образец, патенты на селекционные достижения, свидетельства
        % на программу для электронных вычислительных машин, базу данных, топологию интегральных микросхем,
        % зарегистрированные в установленном порядке.(в ред. Постановления Правительства РФ от 21.04.2016 N 335)
    \end{refsection}%
    \begin{refsection}[bl-author, bl-registered]
        % Это refsection=2.
        % Процитированные здесь работы:
        %  * попадают в авторскую библиографию, при usefootcite==0 и стиле `\insertbiblioauthorimportant`.
        %  * ни на что не влияют в противном случае
        \nocite{vakbib2}%vak
        \nocite{patbib1}%patent
        \nocite{progbib1}%program
        \nocite{bib1}%other
        \nocite{confbib1}%conf
    \end{refsection}%
        %
        % Всё, что вне этих двух refsection, это refsection=0,
        %  * для диссертации - это нормальные ссылки, попадающие в обычную библиографию
        %  * для автореферата:
        %     * при usefootcite==0, ссылка корректно сработает только для источника из `external.bib`. Для своих работ --- напечатает "[0]" (и даже Warning не вылезет).
        %     * при usefootcite==1, ссылка сработает нормально. В авторской библиографии будут только процитированные в refsection=0 работы.
}
